% !TITLE 饮酒·其五
% !AUTHOR 陶渊明
% !DYNASTY 魏晋
% !GENRE 五言古诗
% !ORDER 2

\begin{poem}
结庐\textsuperscript{1}在人境\textsuperscript{2}，而无车马喧。
问君何能尔\textsuperscript{3}？心远地自偏。
采菊东篱下，悠然\textsuperscript{4}见南山。
山气日夕佳，飞鸟相与\textsuperscript{5}还。
此中有真意，欲辨\textsuperscript{6}已忘言。
\end{poem}

\begin{annotations}
\notetext{1}{结庐：建造住宅。庐，简陋的茅草房。}
\notetext{2}{人境：有人居住的世俗环境。}
\notetext{3}{尔：如此，这样。}
\notetext{4}{悠然：安闲舒适、自由自在的样子。这二字点出了诗人无欲无求的洒脱状态。}
\notetext{5}{相与：结伴，共同。}
\notetext{6}{辨：辨析，分辨。这里指想解释清楚其深意。}
\end{annotations}

\begin{appreciation}
《饮酒·其五》是中国文学史上写隐逸最著名的诗篇之一。作者陶渊明，就是《归园田居》里那个辞官种地的陶渊明。他住的地方并不是什么深山老林。

结庐在人境，而无车马喧——屋子就盖在人间，却听不到车马的喧闹。车马是热闹，是应酬，是当官的排场。问君何能尔？心远地自偏——凭什么能做到？心离得远，地方自然就偏了。这是全诗的诗眼：环境吵不吵，不取决于地方，取决于心。心静，闹市也是深山；心躁，深山也睡不着。

采菊东篱下，悠然见南山——在东边的篱笆下采菊花，一抬头，南山就在眼前。这里有个字很讲究：是"见"不是"望"。望是故意去看，见是不经意撞见。陶渊明不是在找南山，采着采着花，抬头撞见了——这种不刻意的遇见，就是"悠然"。

山气日夕佳，飞鸟相与还——傍晚的山色最好，飞鸟结伴归巢。人在采菊，鸟在归巢，都是各自回家。此中有真意，欲辨已忘言——这里面有真正的意思，想说出来，话到嘴边却忘了。"忘言"不是没话说，是不必说。真意这种东西，一说破，就不是它了。

还记得《天问》吗？屈原追着天问了几百个问题，问天，问地，问人为什么命苦。陶渊明反着来：他也看见了那个"真意"，但不问了，说忘了怎么讲。从屈原到陶渊明，中国诗从"问"走到了"忘"。一个把问题抛给天，一个把答案留在心里。

这首诗对你有用的时候，大概是写作业写不进去的时候。心远地自偏——把心收回来，客厅再吵也能当书房；心不在，图书馆也白搭。别信"环境决定一切"。
\end{appreciation}
